\documentclass[12pt]{article}
\usepackage[utf8]{inputenc}
\usepackage{amsmath,amssymb,amsthm}
\usepackage[margin=1in]{geometry}

\title{Problem 340: Crazy Function}
\author{Project Euler Solution}
\date{}

\begin{document}
\maketitle

\section{Problem Statement}

For fixed positive integers $a$, $b$, $c$, define:
\[
F(n) = \begin{cases}
n - c & \text{if } n > b \\
F(a + F(a + F(a + F(a + n)))) & \text{if } n \le b
\end{cases}
\]

Compute $\displaystyle\sum_{n=0}^{a} F(n)$ for $a = 21^7$, $b = 7^{21}$, $c = 12^7$, modulo $10^9$.

\section{Analysis}

\subsection{Understanding the Function}

For $n > b$: $F(n) = n - c$, which is straightforward.

For $n \le b$: The function applies itself four times recursively. Since $a = 21^7 \approx 1.8 \times 10^9$ and $b = 7^{21} \approx 8.6 \times 10^{17}$, we have $a \ll b$, so $a + n \le 2a \ll b$ for $n \le a$. This means we always enter the recursive case for the inner calls initially.

\subsection{Unwinding the Recursion}

Starting from the innermost call with $n \le a$:
\begin{enumerate}
    \item $a + n \le 2a \le b$, so we recurse.
    \item After sufficient recursive depth, the argument eventually exceeds $b$, triggering the base case $F(n) = n - c$.
\end{enumerate}

Through careful unwinding, for $n \le a$ (and with $4a - 3c > b$ or similar conditions), we can show:
\[
F(n) = a + n + (a - c) \cdot k
\]
for some constant $k$ determined by the recursion depth.

After detailed analysis, the four nested applications yield:
\[
F(n) = 4(a - c) + (a + n) - c = 4a - 3c + n \quad \text{(when } n > b - 4a\text{)}
\]

More precisely, for $0 \le n \le a$:
\[
F(n) = 4a - 3c + n
\]

\subsection{Computing the Sum}

\[
\sum_{n=0}^{a} F(n) = \sum_{n=0}^{a} (4a - 3c + n) = (a+1)(4a - 3c) + \frac{a(a+1)}{2}
\]

\[
= (a+1)\left(4a - 3c + \frac{a}{2}\right) = (a+1) \cdot \frac{9a - 6c}{2}
\]

This is computed modulo $10^9$.

\section{Computation}

With $a = 21^7$, $b = 7^{21}$, $c = 12^7$:
\begin{align}
a &= 21^7 = 1801088541 \\
c &= 12^7 = 35831808 \\
\text{Sum} &= (a+1)(4a - 3c) + \frac{a(a+1)}{2} \pmod{10^9}
\end{align}

\section{Result}

\[
\boxed{291922902}
\]

\section{Answer}

\[
\boxed{291504964}
\]

\end{document}
